%% Template for a preprint Letter or Article for submission
%% to the journal Nature.
%% Written by Peter Czoschke, 26 February 2004
%%

\documentclass{nature}
\usepackage{graphicx}
\usepackage{amsmath}


%% make sure you have the nature.cls and naturemag.bst files where
%% LaTeX can find them

\makeatletter
\let\saved@includegraphics\includegraphics
\AtBeginDocument{\let\includegraphics\saved@includegraphics}
\renewenvironment*{figure}{\@float{figure}}{\end@float}
\makeatother

\title{Working title : Collective mind}

%% Notice placement of commas and superscripts and use of &
%% in the author list

\author{Seungwoong Ha$^{1}$, Henrik Olsson$^{1,2}$, Mirta Galesic$^{1,2,3}$}

\begin{document}

\maketitle

\begin{affiliations}
 \item Santa Fe Institute, Santa Fe, NM 87501, USA
 \item Complexity Science Hub Vienna, 1080 Vienna, Austria
 \item Vermont Complex Systems Center, University of Vermont, Burlington, VM 05405, USA
\end{affiliations}

\begin{abstract}

  ABSTRACT

\end{abstract}

\section{Introduction}

INTRODUCTION

\section{Result}

RESULTS

\section{Conclusion and Outlook}

CONCLUSION

\begin{methods}

\subsection{Collective mind model for online news community}

The computational model for the belief dynamics of the online news community consists of three components and their interactions; general topic network where the real-world events are generated, community topic network where articles and comments are posted, and filter between them. At a high level, the model evolves at each time step as follows: (1) real-world events with respective magnitudes are generated according to the general topic network, (2) the filter of each community determines how many of the articles from each event will pass through itself and get posted to their community, and (3) accepted articles and community topic network determines the community topic network of the subsequent time step.

At any given time $t$, the general topic graph from the previous time $t-1$ is expressed as $G^g_{t-1} = (V^g, E^g, F_{t-1}^g, W_{t-1}^g)$, where $V^g$ and $E^g$ denotes the set of vertices and edges, respectively. Here, we assume the topic vertices and edges between them are persistent through time. Each vertex $v_{i}^g \in V^g$ indicates a single topic, and have a normalized frequency value $f_{i,{t-1}}^g \in F_{t-1}^g$ where $F_{t-1}^g$ is the set of all normalized frequency values at time ${t-1}$ and $\sum_i f^g_{i, t-1} = 1$. Each edge $e_{ij,{t-1}}^g \in E_{t-1}^g$ indicates the semantic closeness between two topics, and has a weight value $w_{ij,{t-1}}^g \in W_{t-1}^g$ where $0 \leq w_{ij,{t-1}}^g \leq 1$ and $W_{t-1}^g$ is the set of all weight values at time ${t-1}$. There are total $|V^g| = N$ vertices (topics) and $|E^g| = N(N-1)/2$ edges, since the graph is complete without self-loop. Also, we consider $K$ different community with respective community topic graph $G^{c_k}_{t-1} = (V^{c_k}, E^{c_k}, F_{t-1}^{c_k}, W_{t-1}^{c_k})$ at time ${t-1}$, which shares vertices and edges $(V^{c_k} = V^g, E^{c_k} = E^g)$ but with (potentially) different values for $F_{t-1}^{c_k}$ and $W_{t-1}^{c_k}$.

At each time step $t$, the general topic graph $G^g_t$ generates a new set of events $X_{t}=\{x_{1, t}, \allowbreak  x_{2, t}, \cdots, x_{N_x, t}\}$ for the current timestep, where $N_x$ denotes the number of events per each time step. Each event is consists of $N_{v}=3$ number of event topics from the general topic graph, $x_{j, t+1} = \{v_{q_1}^g, v_{q_2}^g, \cdots, v_{q_{N_{v}}}^g\}$. The first event topic is randomly and uniformly sampled from the topic set $V_{t-1}^g$, and $q_n$-th event topic are sampled with probability proportional to the sum of edge weights with all previously sampled topics, namely, $P(v_{q_n}^g = v_{i}^g) \propto \sum_{r<n} e_{iq_r,{t-1}}^g$. For each event, we define magnitudes with respect to the topic graph. The magnitude of event $x_{j, t}$ for given topic graph $a$ is defined as $M^a(x_{j, t}) = \sum_{i=1}^{N_{v}} f_{q_i,t-1}^a$, where $c$ is a scaling constant.

After the event generation, each event $x_{j, t}$ passes through a filter of each community and generates multiple copies of article $a_{j, t}^{c_k}$. The number of articles generated from event $x_{j, t}$ for community $c_k$ is defined as $\#(a_{j, t}^{c_k}) = M^g(x_{j, t}) \times M^{c_k}(x_{j, t})$, where $M^g(x_{j, t})$ and $M^{c_k}(x_{j, t})$ are magnitude of the event $x_{j, t}$ calculated from the general and $k$-th community topic graph, respectively. As a result, we get a set of all articles $A_{t}^{c_k}$ that contains a total of  $\sum_j \#(a_{j, t}^{c_k})$ articles for each community.


\subsection{Memo}


\subsection{Empirical data}
TBD


\end{methods}



\section{References}
\bibliographystyle{naturemag}
\begin{thebibliography}{10}
  \expandafter\ifx\csname url\endcsname\relax
    \def\url#1{\texttt{#1}}\fi
  \expandafter\ifx\csname urlprefix\endcsname\relax\def\urlprefix{URL }\fi
  \providecommand{\bibinfo}[2]{#2}
  \providecommand{\eprint}[2][]{\url{#2}}
  
  \bibitem{feldman1984cultural}
  \bibinfo{author}{Feldman, M.~W.} \& \bibinfo{author}{Cavalli-Sforza, L.~L.}
  \newblock \bibinfo{title}{Cultural and biological evolutionary processes:
    gene-culture disequilibrium}.
  \newblock \emph{\bibinfo{journal}{Proceedings of the National Academy of
    Sciences}} \textbf{\bibinfo{volume}{81}}, \bibinfo{pages}{1604--1607}
    (\bibinfo{year}{1984}).
  
  \end{thebibliography}
  
%% Here is the endmatter stuff: Supplementary Info, etc.
%% Use \item's to separate, default label is "Acknowledgements"

\begin{addendum}
 \item This research was supported by the XXX.
 \item[Author contribution] S.H. did XXX and M.G and H.O. did YYY.
 \item[Code availability] Simulation code and data files are available at XXX.
 \item[Competing Interests] The authors declare that they have no competing financial interests.
 \item[Correspondence] Correspondence and requests for materials should be addressed to Mirta Galesic (galesic@santafe.edu).
\end{addendum}
\end{document}
